\subsection{THE DIVISION RING OF QUATERNIONS}

It follows from Theorem 1 that the field R is a maximal ordered field, and therefore the only non-commutative division ring of finite rank over R is (up to isomorphism) the division ring of quaternions over R; it is denoted by H and is called the division ring of real quaternions (or simply the division ring of quaternions, when there is no fear of confusion). Since H is of rank 4 over the field R, we can define a topology on H homeomorphic to that of \(R^{4}\) (Chapter VI, § 1, no. 5). To be precise, we shall usually identify H with \(R^{4}\) , the elements 1, i, j, k of the canonical basis of H being identified respectively with the vectors \(e_{0}, e_{1}, e_{2}, e_{3}\) of the canonical basis of \(R^{4}\) .

We recall that the multiplication table of the canonical basis of H is given by the formulae

\[
\begin{array}{c} i ^ {2} = j ^ {2} = k ^ {2} = - \mathrm{I}, \quad i j = - j i = k, \\ j k = - k j = i, \quad k i = - i k = j. \end{array}
\]

The topology of H is compatible not only with the ring structure of H (Chapter VI, § 1, no. 5) but also with its division ring structure; for if x is a non-zero quaternion, the coordinates of \(x^{-1}\) are rational functions of those of x, whose denominators do not vanish. The division ring H, endowed with this topology, is therefore a non-commutative topological division ring. The quaternions \(a + bi\) (a, b real) form a (topological) subfield of H, isomorphic to the field C, with which it is often identified.

We have thus a third example of a locally compact, connected topological division ring, the others being R and C. In fact these are the only topological division rings with these two properties.

We know from algebra that the reduced norm of a quaternion \(x = x_{0} + x_{1}i + x_{2}j + x_{3}k\) is

\[
\mathrm{N} (\boldsymbol {x}) = x _ {0} ^ {2} + x _ {1} ^ {2} + x _ {2} ^ {2} + x _ {3} ^ {2} = | | \boldsymbol {x} | | ^ {2}
\]

(it is therefore the square of the Euclidean norm of x). Since

\[
\mathrm{N} (\boldsymbol {x y}) = \mathrm{N} (\boldsymbol {x}) \mathrm{N} (\boldsymbol {y}),
\]

it follows that the set of all quaternions of norm 1, which is identical with the sphere \(S_{3}\) , forms a compact subgroup of the multiplicative group \(H^{*}\) of non-zero quaternions.

PROPOSITION 2. The multiplicative group \(\mathbf{H}^*\) of non-zero quaternions is a topological group isomorphic to the product of its subgroups \(\mathbf{R}_+^*\) and \(\mathbf{S}_3\) .

Every quaternion \(x \neq 0\) can be written as \(x = ||x||.z\) where \(z\) is a quaternion with norm 1; since \(||xx'|| = ||x||.||x'||\) , the mapping \(x \to (||x||, x/||x||)\) of \(H^*\) onto \(R_+^* \times S_3\) is an isomorphism of the group structures, and from Chapter VI, § 2, no. 3, Proposition 3 it is a homeomorphism of \(H^*\) onto \(R_+^* \times S_3\) .

Remarks. 1) With the use of the relations \(||x + y|| \leqslant ||x|| + ||y||\) and \(||xy|| = ||x||.||y||\) it can be proved directly, as with the field of complex numbers in no. 2, that the topology of \(\mathbf{R}^4\) is compatible with the division ring structure of H (cf.~Chapter IX, § 3, no. 2).

\begin{enumerate}
\def\labelenumi{\arabic{enumi})}
\setcounter{enumi}{1}
\item
  From what has been proved it follows that the spheres \(S_{1}\) and \(S_{3}\) can carry a group structure compatible with their topology. We shall see later that, for each integer n other than 1 and 3, there exists no group structure on \(S_{n}\) compatible with the topology of \(S_{n}\) .
\item
  Every point of the group \(\mathbf{S}_3\) has a neighbourhood homeomorphic to \(\mathbf{R}^3\) (Chapter VI, § 2, no. 4, Proposition 5) but \(\mathbf{S}_3\) is not locally isomorphic to \(\mathbf{R}^3\) ; for if it were it would be abelian, since it is connected (Chapter VII, § 2, no. 2, Theorem 1), and this is not the case, since \(i\) and \(j\) belong to \(\mathbf{S}_3\) and \(ij \neq ji\) (cf.~Chapter V, § 3).
\end{enumerate}

